% =====================================================================================
%  The MOND acceleration scale as dark-sector pressure
%  Carl P. Zimmerman, Briar Creek Tech
%
%  RevTeX 4.2 source.  Compile: pdflatex PAPER_A0_AS_DARK_PRESSURE  (x2, + bibtex if desired)
%  Deliberately framed as a THEORY paper, not a coincidence paper: the coefficient is
%  stated as MEASURED in the abstract, and the prior art is named there too.
% =====================================================================================
\documentclass[aps,prd,twocolumn,superscriptaddress,nofootinbib,floatfix]{revtex4-2}

\usepackage{amsmath,amssymb,graphicx,bm}
\usepackage[colorlinks=true,linkcolor=blue,citecolor=blue,urlcolor=blue]{hyperref}

\newcommand{\azero}{a_{0}}
\newcommand{\Acal}{\mathcal{A}}
\newcommand{\Fcal}{\mathcal{F}}

\begin{document}

\title{The MOND acceleration scale as dark-sector pressure:\\
a derived $\azero(z)$ that switches MOND off at recombination}

\author{Carl P. Zimmerman}
\email{}
\affiliation{Briar Creek Tech, Charlotte, North Carolina, USA}

\date{\today}

\begin{abstract}
In relativistic completions of Milgrom's modified dynamics the acceleration scale $\azero$ is
inserted as a constant of the theory and fitted to galaxy data. We observe that in the Aether
Scalar Tensor theory (AeST) of Skordis and Z\l osnik --- currently the only MOND-like
relativistic theory that reproduces the cosmic microwave background --- the dark sector already
supplies a pressure, and that identifying
$\Acal(Q) \equiv \azero^{2}(Q) = \kappa^{2} G\,[-K(Q)]$
removes $\azero$ as an independent constant. Two consequences follow without further
adjustment. First, $\azero$ is fixed by the dark-energy density, $\azero = \kappa c \sqrt{G
\rho_{\Lambda}}$, giving $9.36\times10^{-11}\,\mathrm{m\,s^{-2}}$ for $\kappa = 1/2$; we
emphasise that $\kappa$ here is \emph{measured}, $\kappa = 0.551 \pm 0.043$, and is
\emph{not} derived --- four candidate values lie within $2\sigma$. Second, and independently of
$\kappa$, the dark-sector pressure evolves, so $\azero$ must evolve with it in a form fixed by
the same equation: $\azero^{2}(z)/\azero^{2}(0) = \sqrt{1+\nu_{0}^{2}}\,/\,\sqrt{1 +
\nu_{0}^{2}(1+z)^{6}}$, which rises monotonically to a maximum today and gives
$\azero(z{=}1090)/\azero(0) = 6.0\times10^{-3}$. MOND is therefore absent when the CMB is
imprinted, making the CMB pass structural rather than fitted. That $\azero$ may depend on the
cosmological scalar-field value, and hence evolve, is not new --- it was computed for TeVeS by
Bekenstein and Sagi, who obtained $\azero$ \emph{decreasing} with time; the identification with
the pressure specifically, the opposite sign of evolution, and the AeST setting are what is
new here. We report what the construction passes ($c_{T}=1$ exactly, no ghost, $\gamma_{\rm
PPN}=1$, a CLASS CMB fit, SPARC rotation curves at $0.108$~dex, the baryonic Tully--Fisher
relation as a theorem, weak lensing from $40$~kpc to $2.2$~Mpc, and the solar system), and we
state plainly what it does not: galaxy clusters remain short by a factor $2.08$ at $R_{500}$,
and we prove that this construction cannot supply the deficit, because the available amplitude
reduces to a charge-abundance ratio that is $10^{2}$--$10^{3}$ too small even if the dark
sector's conserved charge were the entire dark matter content of the universe.
\end{abstract}

\maketitle

% =====================================================================================
\section{Introduction}
% =====================================================================================

Rotation curves of disc galaxies do not fall in the way the visible baryons require. Milgrom's
modified dynamics \cite{Milgrom1983} accounts for this with remarkable economy: given the
baryonic mass distribution it predicts the rotation curve across five decades in mass, and the
radial acceleration relation \cite{McGaugh2016} is tight enough that the scatter is comparable
to the observational uncertainties. The construction contains a single new constant, an
acceleration $\azero \simeq 1.2\times10^{-10}\,\mathrm{m\,s^{-2}}$, above which dynamics is
Newtonian and below which it is not.

The numerical proximity of $\azero$ to accelerations built from cosmological quantities ---
$cH_{0}$, or equivalently $c\sqrt{G\rho_{\Lambda}}$ --- has been remarked on since the earliest
MOND papers and is discussed at length in the review literature \cite{FamaeyMcGaugh2012}. It has
nonetheless remained a coincidence: in every working formulation $\azero$ is a number fitted to
galaxies, and its resemblance to cosmology carries no consequence. Milgrom has proposed specific
normalisations, most recently with coefficient $1/2\pi$ \cite{Milgrom2020}.

Relativistic completions have been harder. Bekenstein's TeVeS \cite{Bekenstein2004} reproduced
the phenomenology but did not survive the cosmic microwave background. The situation changed with
the Aether Scalar Tensor theory (AeST) of Skordis and Z\l osnik \cite{SkordisZlosnik2021}, which
recovers MOND in quasi-static systems \emph{and} fits the CMB and the matter power spectrum. AeST
is the setting for everything below, and the results of this paper are consequences of their
theory rather than of a new one; in their formulation $\azero$ remains an inserted constant.

The observation of this paper is that AeST's dark sector already contains a quantity with the
right dimensions to supply $\azero$, namely its own pressure, and that identifying the two
removes a free constant and forces an evolution law.

% =====================================================================================
\section{The identification}
% =====================================================================================

We take AeST in the form of Ref.~\cite{SkordisZlosnik2021}: a metric $g_{\mu\nu}$, a
unit-timelike vector $A^{\mu}$ with $A^{\mu}A_{\mu} = -1$, and a scalar $\phi$, with the
invariants
\begin{equation}
  Y = q^{\mu\nu}\nabla_{\mu}\phi\nabla_{\nu}\phi , \qquad Q = A^{\mu}\nabla_{\mu}\phi ,
\end{equation}
where $q^{\mu\nu} = g^{\mu\nu} + A^{\mu}A^{\nu}$. The dark sector enters through a function
$\Fcal(Y,Q)$ whose $Q$-dependent part $K(Q)$ carries the pressure, $p = K$, and whose
$Y$-dependent part supplies the MOND limit with a scale $\Acal = \azero^{2}$ inserted by hand.

We instead set
\begin{equation}
  \boxed{\;\Acal(Q) \;=\; \azero^{2}(Q) \;=\; \kappa^{2} G\,[-K(Q)]\;}
  \label{eq:promotion}
\end{equation}
with $\kappa$ dimensionless. In words: the MOND scale is the dark sector's pressure, expressed
as an acceleration. Nothing is added to the theory --- $K(Q)$ was already present, doing
another job --- and one constant is removed.

Two remarks fix the interpretation. Since $w = -1$ requires $p = -\rho$ and the sector's exact
thermodynamics give $n = K'$, $\rho = QK' - K$, $p = K$, the background sits near $K' = 0$; we
return in Sec.~\ref{sec:evolution} to the fact that it sits \emph{near} rather than \emph{at}
that point, which is what makes the evolution law non-trivial. And because
Eq.~(\ref{eq:promotion}) ties $\Acal$ to the same object that drives cosmic acceleration, the
value of $\azero$ becomes a statement about $\rho_{\Lambda}$.

% =====================================================================================
\section{First consequence: the scale is fixed, the coefficient is measured}
% =====================================================================================

With $-K = \rho_{\Lambda}c^{2}$ on the background, Eq.~(\ref{eq:promotion}) gives
\begin{equation}
  \azero \;=\; \kappa\, c\, \sqrt{G \rho_{\Lambda}} \;=\; \frac{c H_{\Lambda}}{Z},
  \qquad Z \equiv \frac{\sqrt{8\pi/3}}{\kappa},
  \label{eq:a0}
\end{equation}
the second form following because the Friedmann equation supplies $H_{\Lambda} =
\sqrt{8\pi/3}\,\sqrt{G\rho_{\Lambda}}$. For $\kappa = 1/2$, equivalently $Z = 2\sqrt{8\pi/3} =
\sqrt{32\pi/3} = 5.78881$, this yields
\begin{equation}
  \azero = 9.3619\times 10^{-11}\ \mathrm{m\,s^{-2}} .
\end{equation}
Using $\rho_{\rm total}$ and $cH_{0}$ in place of $\rho_{\Lambda}$ and $cH_{\Lambda}$ --- a
defensible alternative bookkeeping which we carry throughout as the ``alt'' footing --- gives
$1.1279\times10^{-10}\,\mathrm{m\,s^{-2}}$ instead. We quote both for every dimensionful result.

\paragraph*{The coefficient is not derived.} We state this in the strongest terms because the
literature is full of near-misses. A distance-free determination from the radial acceleration
relation gives $\kappa = 0.551 \pm 0.043$, and the baryonic Tully--Fisher route gives $0.465 \pm
0.076$; $\kappa = 1/2$ lies $1.19\sigma$ from the former, and at least four candidate values sit
within $2\sigma$ of one another. The data are therefore consistent with $\kappa = 1/2$ but do not
select it.

Furthermore, the apparent depth of the question is illusory. Because the $\pi$ in
Eq.~(\ref{eq:a0}) cancels between $\sqrt{8\pi/3}$ and $Z$, the statement $\kappa = 1/2$ is
identically the statement $Z = 2\sqrt{8\pi/3}$, i.e.\ that $Z$ is exactly twice the Friedmann
factor. Any normalisation argument that returns $\kappa^{2} = 8\pi\epsilon$ and then reports
$\epsilon = 1/32\pi$ has not derived anything: the $8\pi$ is carried over from its own left-hand
side, and the residual $4$ in $32\pi$ is $\kappa^{-2}$. The entire open content is one factor of
two, and we do not claim to have it.

By contrast, the coefficient is not what the next section rests on.

% =====================================================================================
\section{Second consequence: a derived $\azero(z)$}
\label{sec:evolution}
% =====================================================================================

The dark-sector pressure evolves, so Eq.~(\ref{eq:promotion}) forces $\azero$ to evolve. Writing
$u = Q - Q_{0}$ and $\nu = \mu_{p}u/M^{2}$ for the sector's dimensionless displacement, with the
brane-form $K(u) = -M^{4}\sqrt{1 - \mu_{p}^{2}u^{2}/M^{4}}$ at $\beta \equiv
\mu_{p}^{2}\Lambda_{D}^{2}/M^{4} = 1$, and using $n \propto a^{-3}$ so that $\nu(z) =
\nu_{0}(1+z)^{3}$, Eq.~(\ref{eq:promotion}) gives
\begin{equation}
  \boxed{\;\frac{\azero^{2}(z)}{\azero^{2}(0)}
  \;=\; \frac{\sqrt{1+\nu_{0}^{2}}}{\sqrt{1+\nu_{0}^{2}(1+z)^{6}}}\;}
  \label{eq:a0z}
\end{equation}
with a single parameter $\nu_{0}$. Constraints already carried by the sector confine $\nu_{0}
\in [2.14, 17.7]\times10^{-5}$, whence a transition redshift $z_{t} = \nu_{0}^{-1/3} - 1 \in
[17,35]$ and, at recombination,
\begin{equation}
  \azero(z{=}1090)\,/\,\azero(0) \;=\; 6.0\times10^{-3}.
\end{equation}

Three features deserve emphasis. (i) Eq.~(\ref{eq:a0z}) is \emph{monotonically rising} toward
the present: $\azero$ is at its maximum today. (ii) It is exactly $\kappa$-independent, so it
tests the identification (\ref{eq:promotion}) without depending on the unresolved coefficient.
(iii) MOND is therefore \emph{off} when the CMB is imprinted. This last point is the reason the
CMB is not a struggle here: the clustering behaviour required at $z \sim 10^{3}$ is a consequence
of Eq.~(\ref{eq:a0z}) rather than an assumption, and the failure mode that removed TeVeS from
contention does not arise.

Note also that $\nu_{0} \neq 0$ is required for Eq.~(\ref{eq:a0z}) to be non-empty, so the
background sits at $u_{0} = \nu_{0}\Lambda_{D} \neq 0$ and $w = -1$ holds only to
$\mathcal{O}(\nu_{0}^{2})$, i.e.\ to a fractional offset in $[4.6\times10^{-10},
3.1\times10^{-8}]$. That is far below any current determination of $w_{0}$, but it should be
written as $w = -1 + \mathcal{O}(\nu_{0}^{2})$ rather than as exactly $-1$.

\paragraph*{Relation to prior work.} That $\azero$ may depend on the cosmological value of the
dark-sector scalar, and hence evolve, is \emph{not} new. Bekenstein and Sagi computed exactly
this for TeVeS \cite{BekensteinSagi2008}, finding $\azero$ varying as the exponential of the
scalar field and \emph{decreasing} with time; Ref.~\cite{FamaeyMcGaugh2012} treats the general
statement --- that $\azero$ for a quasi-static system depends on the field value at the epoch the
system collapsed --- as established background. What is specific to Eq.~(\ref{eq:promotion}) is
the identification with the \emph{pressure} rather than with an exponential of the field, so that
no new function is introduced; the \emph{opposite sign} of the resulting evolution; and the AeST
setting, which is the one that survives the CMB. We make no priority claim beyond that, and we
have not performed an exhaustive literature search.

% =====================================================================================
\section{What the construction passes}
% =====================================================================================

We summarise, with the caveat that all fitted numbers below use the algebraic interpolation
$g^{2} = g_{\rm bar}^{2} + \azero g_{\rm bar}$, equivalently Milgrom's $\nu(y) = \sqrt{1+1/y}$
\cite{Milgrom1999}, unless stated otherwise.

\begin{itemize}\itemsep2pt
\item \textbf{Tensor sector.} $c_{T} = 1$ exactly, as in the parent theory, so GW170817 is
      satisfied identically. The no-ghost condition of the parent theory transfers: with the
      promotion the two-field longitudinal conditions reduce to the two standard AQUAL
      conditions.
\item \textbf{Lensing.} $\gamma_{\rm PPN} = 1$. A fit to the weak-lensing radial acceleration
      relation of Ref.~\cite{Mistele2024} from $40$~kpc to $2.2$~Mpc, with no free amplitude,
      gives $\chi^{2}/\mathrm{dof} = 2.03$ (canonical footing) and $0.94$ (alt).
\item \textbf{CMB.} A CLASS run with Eq.~(\ref{eq:a0z}) in place returns $\Delta\chi^{2} = 1.34$
      over $4998$ multipoles relative to the unmodified fit, with $\theta_{*}$ displaced by
      $0.18$ of the Planck uncertainty.
\item \textbf{Galaxies.} SPARC rotation curves at $0.108$~dex scatter with stellar
      mass-to-light ratio $\Upsilon = 0.70$. The baryonic Tully--Fisher relation $v^{4} = G
      M_{b}\azero$ follows as a \emph{theorem} of the interpolation rather than as a fit, since
      $r$ drops out identically, with $\mathrm{d}\ln v/\mathrm{d}\ln \azero = 1/4$.
\item \textbf{Solar system.} The algebraic interpolation above approaches Newtonian dynamics as
      $g_{\rm bar} + \azero/2$, i.e.\ with a constant residual anomaly, which exceeds planetary
      ranging bounds by a factor $\sim 1.3\times10^{3}$. It must therefore be read as an
      effective description valid where the data lie, with the exponential form $\nu(y) = [1 -
      e^{-\sqrt{y}}]^{-1}$ of Ref.~\cite{McGaugh2016} as the interpolation proper; that form
      approaches Newtonian exponentially and the anomaly at Earth is unmeasurably small. The two
      differ by at most $0.057$~dex over $g_{\rm bar}/\azero \in [0.03,10]$, i.e.\ half the SPARC
      scatter, so the fits above are unaffected; but every quoted number is
      interpolation-conditional and we flag it as such.
\end{itemize}

% =====================================================================================
\section{What it does not pass, and a no-go for repairing it}
% =====================================================================================

\paragraph*{Galaxy clusters.} Taking the baryons of the twelve X-COP clusters
\cite{Eckert2019,Ettori2019} and predicting the dynamical mass from the interpolation above, the
median discrepancy is a factor $2.08$ at $R_{500}$ (canonical footing; $1.92$ alt), rising to
$4.21$ inside $0.1R_{500}$. This is the long-standing cluster problem of MOND-like theories
\cite{FamaeyMcGaugh2012}; it is not specific to Eq.~(\ref{eq:promotion}) and it is not solved by
AeST as published.

\paragraph*{And Eq.~(\ref{eq:promotion}) cannot repair it.} The natural repair is to let the
local dark-sector charge modify the effective interpolation, since $\Acal$ is now
$Q$-dependent. Expanding, and using $\Acal' = \kappa^{2}G(-K') = -\kappa^{2}G n$ together with
$-K = \rho_{\Lambda}c^{2}$ and $\rho_{\rm charge} = Q_{0}n$, the induced correction is
\begin{equation}
  \Delta\mu \;=\; -\,\frac{\rho_{\rm charge}}{\rho_{\Lambda}c^{2}}\;\Phi\;T(y),
  \qquad T = y\,\Fcal''(y),
  \label{eq:cluster}
\end{equation}
in which $\nu_{0}$, $\Lambda_{D}$, $Q_{0}$, $M$ and $\kappa$ all cancel: the amplitude is a pure
charge-abundance ratio. Reproducing the cluster deficit requires a coefficient of order
$0.5$; Eq.~(\ref{eq:cluster}) supplies at most $\sim 10^{-9}$. Even setting the dark sector's
conserved charge to the \emph{entire} dark matter content, $\Omega_{\rm kd} = \Omega_{\rm dm}$,
leaves the term short by a factor $10^{2}$--$10^{3}$; matching would require $\Omega_{\rm kd}
\sim 2.7\times10^{2}$. The charge is a trace species, and it is one for the same reason
Eq.~(\ref{eq:a0z}) switches MOND off early. We regard this as a genuine no-go for this
particular repair, not as a bound to be improved.

\paragraph*{Two further open problems.} The sector's pressureless excitation, which the CMB fit
requires at the level of $\Omega_{\rm dm}$, is forced cold by that same fit and has no evident
mechanism halting its collapse in galaxies; followed through, the endpoint is inconsistent with
the observed mass of Sgr~A$^{*}$ by some five orders of magnitude. And $\kappa$ remains measured
rather than derived, as stated in Sec.~III. We list these rather than defer them.

% =====================================================================================
\section{Falsifiable predictions}
% =====================================================================================

\begin{enumerate}\itemsep2pt
\item \textbf{No $\azero$ evolution below $z \sim 5$.} Eq.~(\ref{eq:a0z}) is flat to
      $\sim10^{-8}$ at low redshift. Any robust detection of $\azero$ evolution below $z\sim5$
      falsifies it. Current intermediate-redshift determinations of the radial acceleration
      relation are in mild tension with this and constitute the sharpest near-term test.
\item \textbf{Suppressed rotation velocities at cosmic dawn.} Since $v^{4} = GM_{b}\azero$ is a
      theorem, $v(z)/v(0) = [\azero(z)/\azero(0)]^{1/4}$, giving deficits of $0.34\%$, $7.8\%$
      and $18.9\%$ at $z = 10$, $20$ and $30$. The effect becomes detectable near $z \approx 17$.
\item \textbf{Wide binaries.} The framework predicts a velocity-boost ratio $\gamma_{v} =
      1.2139$ (canonical) or $1.2592$ (alt) in the Gaia DR4 wide-binary sample, against $1$ for
      Newtonian dynamics; the two footings are separated by $2.68\sigma$, so DR4 distinguishes
      them. This prediction was registered before the data.
\end{enumerate}

% =====================================================================================
\section{Conclusion}
% =====================================================================================

Identifying the MOND acceleration scale with the dark sector's pressure,
Eq.~(\ref{eq:promotion}), removes $\azero$ as an independent constant of AeST and forces an
evolution law, Eq.~(\ref{eq:a0z}), which is $\kappa$-independent, rises to a maximum today, and
puts $\azero$ at $0.6\%$ of its present value at recombination. The CMB pass of the parent theory
thereby becomes structural. The construction retains $c_{T}=1$, no ghost, $\gamma_{\rm PPN} = 1$,
the SPARC relation at $0.108$~dex, the Tully--Fisher relation as a theorem, and weak lensing to
$2.2$~Mpc. It does not solve galaxy clusters, and we have shown that it cannot: the available
amplitude is a charge-abundance ratio too small by $10^{2}$--$10^{3}$. The coefficient $\kappa$
is measured, not derived, and the open content of that question is a single factor of two.

\begin{acknowledgments}
This work is built entirely on the Aether Scalar Tensor theory of Skordis and Z\l osnik, and on
the empirical relations assembled by McGaugh, Lelli and Schombert and by the X-COP
collaboration. All numerical claims are reproducible from scripts in the accompanying public
repository.
\end{acknowledgments}

% =====================================================================================
\begin{thebibliography}{99}

\bibitem{Milgrom1983} M.~Milgrom, Astrophys.\ J.\ \textbf{270}, 365 (1983).
\bibitem{McGaugh2016} S.~S.~McGaugh, F.~Lelli, and J.~M.~Schombert,
  Phys.\ Rev.\ Lett.\ \textbf{117}, 201101 (2016).
\bibitem{FamaeyMcGaugh2012} B.~Famaey and S.~S.~McGaugh,
  Living Rev.\ Relativity \textbf{15}, 10 (2012).
\bibitem{Milgrom2020} M.~Milgrom, Phys.\ Rev.\ D \textbf{102}, 084010 (2020).
\bibitem{Bekenstein2004} J.~D.~Bekenstein, Phys.\ Rev.\ D \textbf{70}, 083509 (2004).
\bibitem{SkordisZlosnik2021} C.~Skordis and T.~Z\l osnik,
  Phys.\ Rev.\ Lett.\ \textbf{127}, 161302 (2021).
\bibitem{BekensteinSagi2008} J.~D.~Bekenstein and E.~Sagi,
  Phys.\ Rev.\ D \textbf{77}, 103512 (2008), arXiv:0802.1526.
\bibitem{Milgrom1999} M.~Milgrom, Phys.\ Lett.\ A \textbf{253}, 273 (1999).
\bibitem{Mistele2024} T.~Mistele, S.~McGaugh, F.~Lelli, J.~Schombert, and P.~Li,
  Astrophys.\ J.\ Lett.\ \textbf{969}, L3 (2024).
\bibitem{Eckert2019} D.~Eckert \textit{et al.}, Astron.\ Astrophys.\ \textbf{621}, A40 (2019).
\bibitem{Ettori2019} S.~Ettori \textit{et al.}, Astron.\ Astrophys.\ \textbf{621}, A39 (2019).

\end{thebibliography}

\end{document}
